\documentclass[10pt,a4paper]{article}
\usepackage[margin=1.7cm]{geometry}
\usepackage[T1]{fontenc}
\usepackage{lmodern}
\usepackage{microtype}
\usepackage{xcolor}
\usepackage{tabularx}
\usepackage{booktabs}
\usepackage{enumitem}
\usepackage{titlesec}
\usepackage{fancyhdr}
\usepackage{tcolorbox}
\usepackage{needspace}
\usepackage[hidelinks]{hyperref}

\definecolor{navy}{HTML}{1F2A44}
\definecolor{coral}{HTML}{E4572E}
\definecolor{soft}{HTML}{F3F4F7}

\titleformat{\section}{\large\bfseries\color{navy}}{\thesection}{0.6em}{}[\vspace{-0.4em}{\color{coral}\rule{\linewidth}{1pt}}]
\titleformat{\subsection}{\normalsize\bfseries\color{navy}}{\thesubsection}{0.5em}{}
\titlespacing*{\section}{0pt}{1.0em}{0.5em}
\titlespacing*{\subsection}{0pt}{0.7em}{0.25em}
\setlist{nosep,leftmargin=1.3em,topsep=2pt}
\setlength{\parindent}{0pt}
\setlength{\parskip}{3pt}
\renewcommand{\arraystretch}{1.15}

\pagestyle{fancy}
\fancyhf{}
\lhead{\small\color{navy}ThemeSync: defense cheat sheet}
\rhead{\small\color{navy}\thepage}
\renewcommand{\headrulewidth}{0.4pt}

\newtcolorbox{key}[1][]{colback=soft,colframe=navy,boxrule=0.6pt,arc=2pt,left=6pt,right=6pt,top=4pt,bottom=4pt,#1}
\newtcolorbox{trap}{colback=coral!7,colframe=coral,boxrule=0.6pt,arc=2pt,left=6pt,right=6pt,top=3pt,bottom=3pt}
\newcommand{\Q}[1]{\par\smallskip\textbf{\color{navy}Q: #1}\par}
\newcommand{\A}{\textbf{A:} }

\begin{document}

\begin{center}
{\LARGE\bfseries\color{navy} ThemeSync: Defense Cheat Sheet}\\[2pt]
{\small Hermoso, Quite, Tercero \;$\cdot$\; based on Manu 0411.3 (printed) and the system on \texttt{main} \;$\cdot$\; 2026-11-04}
\end{center}

\begin{key}
\textbf{How to use this.} Don't memorize sentences. Know the \emph{beats}, the \emph{reason for each method}, and the \emph{numbers} in Section 7. Answer in short sentences: one point each. If you don't know, say what you do know and what the study did not measure.
\end{key}

%=====================================================
\section{The panel and what they will check}

\begin{tabularx}{\linewidth}{@{}p{3.4cm}p{4.6cm}X@{}}
\toprule
\textbf{Panelist} & \textbf{Background} & \textbf{Likely focus} \\
\midrule
Mr.\ Joseph Jessie S.\ Oñate, MSc (Chairman) & AI, NLP, remote sensing, data viz; heads the AI Research Center; deep-learning segmentation (RAU-Net) & YOLO and GNN training, the 200-image dataset, overfitting, boxes vs.\ segmentation masks, the CGL domain gap, why keep the GNN \\
Dr.\ Rosel O.\ Onesa (Member 1) & CCS OIC Dean; data science and NLP & Statistics (Wilcoxon, Holm, Kendall's $W$), sample sizes, purposive sampling, data quality, choice of $k$, the OCR \\
Mr.\ Allan O.\ Ibo Jr., MSc (Member 2) & CSPC Web and Social Media Administrator & Real use by a pubmat team, output formats and print, Canva comparison, where it runs, the dashboard's ease of use \\
\bottomrule
\end{tabularx}

\smallskip
\textbf{They will check their own May comments} (RSC sheet). Have the page ready:
\begin{itemize}
\item Sizes and font in the Visual DNA: p.45 and Figure 2 (p.56). \quad Pomelli: p.2, ref.\ [7] on p.13.
\item How the Visual DNA's correctness is ensured: \emph{Validation of Hybrid Visual DNA Correctness}, p.45--46 (elbow method; JSON Schema).
\item Training data and a sample dataset: p.49--52 and Appendix C (p.112).
\item Evaluation by experts: Phase 4 (p.60), MOS (p.64), Statistical Treatment (p.65), results p.93--99.
\item Correct evaluation metrics: IoU limited to Phase 1 (p.62); F1 defined by constraint satisfaction (p.63).
\end{itemize}
Our sample signatories in the dashboard include ``Dr.\ Rosel O.\ Onesa, Dean, College of Computer Studies''. Expect it to be noticed.

%=====================================================
\section{The question you cannot avoid: the experts preferred the baseline}

\textbf{Expect:} ``Your own experts rated the version without your method higher. What did the thesis achieve?''

\begin{key}
\begin{enumerate}
\item \textbf{We were rigorous.} The statistical plan was fixed and committed (\texttt{b9a60ab}) before the scores were imported, and we report the result as it came out.
\item \textbf{What it means.} Following the poster's structure does not, by itself, make a better certificate or ID. Experts judge the target format on its own, and they penalized uneven gaps. This is Conclusion 3, and it is a valid finding.
\item \textbf{Why our metric missed it.} Formula 3.11 scores collision, alignment, the \emph{amount} of whitespace and legibility, but not how \emph{evenly} the gaps are spread. It still agreed in direction: the baseline also had a slightly better fitness (0.0855 vs.\ 0.0879, lower is better). Post hoc: mapping + GNN left a larger largest gap in 23 of 32 pairs.
\item \textbf{What still holds.} SO1 and SO2 were met: the structure transfers, all 64 rated layouts are valid, edits are never lost, and re-solves are real-time. The designer can fix the spacing in a few edits, and the system keeps them: that is the ``stateful refinement'' in the title.
\end{enumerate}
\end{key}

\begin{trap}
\textbf{``So remove the GNN?''} Don't agree. The GNN does reduce structure error (0.050$\to$0.036 held out; 0.117$\to$0.045 on IDs). The study tested mapping and GNN \emph{together}, so it cannot separate them; Recommendation 3 compares mapping alone, mapping + GNN, and the baseline.\\[2pt]
\textbf{``Have you fixed it?''} An even-gap step exists on a separate branch, built after the evaluation. It has not been rated, so we claim no result. Don't present it as an improvement.\\[2pt]
\textbf{``Was SO3 met?''} SO3 was \emph{to evaluate}. The evaluation was carried out fully as planned; its result did not support the mapping + GNN condition. ``Carried out'' (objective) is not the same as ``supported'' (result).
\end{trap}

%=====================================================
\section{Part 1: the one-minute story}

\textbf{Four beats:} \textbf{Problem $\to$ Theme $\to$ Sync $\to$ Results.}
\begin{itemize}
\item \textbf{Problem:} materials redesigned by hand; a theme change can undo manual edits.
\item \textbf{Theme:} poster in $\to$ Visual DNA $\to$ certificate or event ID out.
\item \textbf{Sync:} you move one element; the rest re-arrange around it, and earlier edits stay.
\item \textbf{Results:} valid, fast; next to fix is the spacing.
\end{itemize}

\begin{key}
\small
Student organizations make the same materials again and again: a poster for an event, then the certificates and IDs that go with it, each redesigned by hand. And in most tools, a change to the theme can undo the designer's manual adjustments.

ThemeSync is a graph-based design orchestration framework that addresses this in two parts. \textbf{Theme:} it reads the organization's master poster (its colours, logos, fonts, and the arrangement of its elements) and builds a Visual DNA from it. That DNA generates new layouts: certificates or event IDs. \textbf{Sync:} every element is a node in a graph. When the designer moves or resizes one, it becomes a Hard constraint, and a constraint solver re-arranges every Soft element around it in real time, without losing earlier edits. This repeats until the designer is satisfied.

We designed, built, and evaluated it. Every layout stayed valid, and the solver recalculated each edit within a tenth of a second. Four expert designers' ratings showed us that even spacing is what we must improve next.
\end{key}

\begin{trap}
``Within a tenth of a second'' is the \emph{solver}: every recalculation (median 0.5\,ms, max 2.1\,ms). A \emph{full} edit (re-solve, re-score, re-draw, save) has a median of 91\,ms, all within 1\,s; only 47\% of re-solving edits finished within 0.1\,s. Never say ``every edit takes under a tenth of a second''.
\end{trap}

%=====================================================
\section{Part 2: the system, stage by stage}

\textbf{The chain:} Poster $\to$ colours, details, elements $\to$ \textbf{Visual DNA} $\to$ structure mapping + GNN $\to$ the user's content $\to$ \textbf{Cassowary} $\to$ \textbf{6 candidates} $\to$ \textbf{Formula 3.11} $\to$ render $\to$ designer edits $\to$ \textbf{state machine} $\to$ re-solve $\to$ repeat.

\subsection*{Step 1: reading the master poster (Phase 1)}
\textbf{1. Palette: K-Means in CIE L*a*b* (F3.1).} Groups pixels into primary, secondary and accent colours. \emph{Why Lab:} distance in Lab matches perceived difference; RGB doesn't. \emph{Why K-Means:} directly minimizes within-cluster spread (WCSS); simple; seeded so the same poster always gives the same palette. $k$ = 2 to 6, best of five seeded starts, chosen at the elbow (drawn in the Step 3 report). This answers the May ``correctness'' comment (p.45).

\textbf{2. Stylistic markers: SIFT (F3.2).} Keypoints (corners, dense texture) measure where the poster is visually busy: (a) a saliency map of the text-free background for the legibility penalty; (b) ranking the motifs YOLO found; (c) choosing artwork to carry over. \emph{Why SIFT:} scale-invariant, so detail density compares fairly across poster sizes; no training; deterministic. \emph{Why not SURF:} patented, outside OpenCV's standard build; SIFT covers the need.

\textbf{3. Elements: YOLOv8s + Tesseract OCR.} Boxes for logo, text, underlay, embellishment; OCR labels which text box is title, header or details. \emph{Why YOLOv8s:} one-stage, fast, fits an RTX 3060, and it outputs \emph{boxes}, which is exactly what the solver needs (x, y, w, h). \emph{1024\,px:} small logos vanish at 640. \emph{Own 200 images:} the CGL-only detector made \textbf{zero} detections at 0.25 on real posters (domain gap: CGL is e-commerce). mAP50 0.728; mean IoU 0.644; recall 75\%. Logos strongest (IoU 0.844), embellishments weakest (0.445). Better annotations helped more than a bigger model.

\textbf{4. Background (Stage 0c).} Text removed; the background rebuilt from the poster's borders (inpainting); title artwork and motifs kept as separate assets.

\textbf{Hybrid Visual DNA} = extracted (palette, saliency, structure, markers) + user-entered (font, weight, size). The user's part is validated by JSON Schema before anything runs: weights 100--900, sizes 8\,px to 15\% of the shorter side, every field present and within its limit (p.46).

\subsection*{Step 2: building the layout (Phases 2--3)}
\textbf{5. Structure mapping.} Target elements with a counterpart on the poster (logo row, header, title, details) are pinned at the poster's positions; the rest go in reading order between them. This is what makes the certificate look like it came from the same poster.

\textbf{6. GNN (F3.6--3.8).} Refines the unmapped elements. Node = [x, y, w, h, colour, saliency]; edge = [$\Delta$x, $\Delta$y, IoU, alignment]. Message passing: an MLP computes messages, they are aggregated, a GRU updates each node. Trained on CGL from deliberately distorted starts; output registered back to the poster's anchors. \emph{Why a graph network:} a layout is a set of elements and relations with a varying count; a GNN handles any number of nodes and is order-invariant. Held-out position error 0.050$\to$0.036.

\textbf{7. Node instantiation.} The content payload becomes nodes: a \emph{Spatial Hierarchy Index} (reading order, 0 = logo row \ldots\ 8 = signatures) and a baseline size per role. CGL labels every text as ``text'', so hierarchy is carried by size and order (``implicit geometric hierarchy'').

\textbf{8. Cassowary (F3.3--3.5).} Relations become linear equations and inequalities, solved for every x, y, w, h. Hard = required (never broken); Soft = weighted (violation minimized). \emph{Why:} built for interactive layout (Apple's Auto Layout is based on it), native required/weighted constraints, incremental re-solve. Median 0.5\,ms.

\textbf{9. Six candidates + Formula 3.11.} Vary archetype, margin, vertical rhythm, motif arrangement. The poster's composition picks the archetype; Formula 3.11 ranks within it. Penalties: collision (overlap area), alignment (distance to grid), whitespace (deviation from 0.65), legibility (saliency under text). Weights (code, not printed): $\alpha$ = 50, $\beta$ = 2, $\gamma$ = 1, $\delta$ = 1.5.

\textbf{10. Rendering.} SVG with palette, fonts, rebuilt background, motifs and title artwork; PNG/PDF via cairosvg.

\textbf{11. State machine (F3.9--3.10).} Only a \emph{finished} action is an event (drop, end of resize, lock, text edit); drag streams are ignored. The edited element becomes Hard, Cassowary re-solves the rest, a new state is appended to the JSON chain. Undo appends an older state; nothing is erased.

\subsection*{Step 3: the Evaluation Report}
Formula 3.11 term by term, F1 per element, manual edits kept, latency, the elbow graph, schema checks; plus the mail-merge download.

\begin{trap}
SIFT does \textbf{not} find the logos; YOLO does. The detector is COCO-pretrained, \textbf{not} CGL-pretrained; CGL trains only the GNN. OCR words are \textbf{never} placed in the output; content comes from the user's validated payload. The GNN's gain after anchoring is small (0.075 vs.\ 0.081 for interpolation); colour and saliency are per-layout, not per-element (stated limitation).
\end{trap}

\textbf{Cross-cutting questions.}
\Q{Where exactly is the graph?}
\A Twice: in the GNN (elements = nodes, relations = edges) and in the solver's constraints, which follow the same relations.
\Q{Why not generative AI?}
\A It produces pixels: flat, not editable, different each run. Administrative documents need exact names, editable vector elements, and the same output for the same input.
\Q{What is ``real-time''?}
\A Nielsen's limits: 0.1\,s feels instant (each solver recalculation) and 1.0\,s keeps the task uninterrupted (each full edit). Both were met.

%=====================================================
\section{Part 3: methodology and evaluation per objective}

\textbf{Research design: Constructive.} The problem required \emph{building} an artifact and then evaluating it, not observing existing tools. Phase 4 is a blind, paired comparison inside that design.

\textbf{Four phases:} (1) extraction and graph construction; (2) CSP and state synchronization; (3) heuristic testing, 96 runs: 16 posters $\times$ 2 formats $\times$ 3 conditions; (4) expert validation, 64 layouts.

\subsection{SO1: Relational Logic and Feature Extraction model}
\textbf{Data:} 40-image validation split (607 elements); 2,900 held-out CGL layouts (a seeded 5\%); 16 unseen posters, three conditions (mapping + GNN, mapping only, baseline).\\
\textbf{Metrics:} IoU, mAP50, recall for the detector (Phase 1 only); position error for the GNN; structure error for the transfer.\\
\textbf{Results:} detector as above; GNN 0.050$\to$0.036; structure error 0.117$\to$0.045 (IDs), 0.089$\to$0.056 (certificates); mapping + GNN equals mapping only, so the transfer comes from the mapping, and the GNN moved unmapped elements by at most 7.5\% of the height.\\
\textbf{Verdict:} met, with caveats: in-domain data was needed; artwork and motifs still need human review; the GNN changed final layouts little.

\Q{Why is IoU used only in Phase 1?}
\A A generated layout in a new format has no ground-truth positions, so there is nothing to overlap with.
\Q{Is 200 images enough? Overfitting?}
\A A held-out 40-image validation split; mAP50 rose from 0.486 to 0.728 when annotations improved. It is small; Recommendation 1 expands it. COCO pretraining helps with little data.
\Q{Why boxes, not segmentation?}
\A The solver works with boxes; a mask would be reduced to a box anyway. Boxes also match the CGL four-class schema.
\Q{How were annotations kept consistent?}
\A Written class definitions; two annotators, the group leader reviewed; an outside designer checked 15 images (boxes correct on all 15, classes on 10), then the set was expanded to 200 and re-annotated (p.51).

\subsection{SO2: Reactive Layout Synthesis engine}
\textbf{Data:} the 64 rated layouts; the synthetic stress test; 34 sessions $\times$ 40 edits = 1,360 edits on the RTX 3060 workstation.\\
\textbf{Metrics:} F1 defined by constraint satisfaction (TP: generated and violating nothing; FP: overlaps or leaves the canvas; FN: missing); latency against 0.1\,s and 1.0\,s.\\
\textbf{Results:} F1 = 1.0 on all 64; every edit kept through later re-solves; solver 100\% within 0.1\,s; full edits 100\% within 1\,s (61\% within 0.1\,s). Stress test: ID header overflows past 48 characters per line; certificate signature blocks overlap from six.\\
\textbf{Verdict:} met; limits are text fit and the time of re-solving edits.

\Q{F1 is 1.0 everywhere. Isn't that a weak metric?}
\A It checks structural validity, which Cassowary guarantees. It can't see text overflow, so we added a text-fit check (Tables 10--11), and Recommendation 2 measures text fit with the F1.
\Q{Why are re-solves slower than the solver?}
\A Re-scoring legibility against the poster's background takes most of the time; on the synthetic theme a full edit took about 4\,ms. Recommendation 2: cheaper legibility scoring.

\subsection{SO3: Heuristic Fitness Function and expert validation}
\textbf{Data:} Phase 3 battery (96 runs, RTX A5000 server); Phase 4 ratings.\\
\textbf{Fitness:} baseline slightly better (p = 0.030). The fitness alone could not pick archetypes (top candidates tied), so the poster's composition picks the archetype and the fitness ranks within it.\\
\textbf{Calibration (Table 15):} rule fixed in advance: keep a weight change only if the selected layouts improve on one design rule and worsen on none. No change was kept; halving or doubling any weight changed $\le$1 of 64. The weights are stable; improving the function needs new \emph{terms}, not new weights.\\
\textbf{Phase 4 design:} 4 professional designers; 16 posters by purposive sampling (real school posters; unseen, verified by an image-duplicate test; text fits both formats) $\times$ 2 formats $\times$ 2 conditions = 64 layouts $\times$ 3 parameters = 768 ratings. Blind: random codes, page order randomized per expert, pages from the same poster never adjacent; printed packets with the master beside each layout; 5-point scale adapted from ITU-T P.800.

\subsection*{The statistics, in plain words}
\begin{itemize}
\item \textbf{MOS} = mean rating (with SD), per parameter and overall (the mean of the three). Overall 3.53, between Acceptable and Good.
\item \textbf{Kendall's $W$}: do the experts \emph{order} the layouts the same way? 0 = no, 1 = perfectly. Overall 0.630 = moderate (Schmidt: 0.5--0.7). ICC(2,1) 0.16--0.31 and Krippendorff's $\alpha$ 0.01--0.23 are low: they ranked alike but used the scale differently (expert means 2.88 to 4.32). So the paired comparison is more dependable than absolute MOS values.
\item \textbf{Wilcoxon signed-rank (two-sided):} ratings are ordinal; the layouts are matched pairs (same poster, same format); 32 pairs is small; no normality assumed. Each layout's score = mean of four experts (the layout is the unit, so the 768 ratings are not treated as independent). Two-sided: no direction assumed, which mattered, since the result went the opposite way.
\item \textbf{Holm correction:} four tests (overall + three parameters) inflate false positives; Holm controls that and is less conservative than Bonferroni. Overall was named the main test in advance.
\item \textbf{Rank-biserial $r$:} effect size from $-1$ to $+1$; negative favours the baseline; $|r| \ge 0.5$ is large. Overall $r = -0.95$: 25 pairs favoured the baseline, 2 mapping + GNN, 5 equal.
\item \textbf{Sensitivity:} paired t-test and Bonferroni agreed. Without the one identical pair, Hierarchy became borderline significant (p = 0.046).
\item \textbf{Post hoc (exploratory, uncorrected):} largest gap 0.072 vs.\ 0.020; gap CV 0.47 vs.\ 0.00; pairs whose gaps grew lost more MOS (Spearman $\rho$ $-0.42$ to $-0.45$). Correlational: it does not prove the cause.
\end{itemize}

\Q{Why only four experts?}
\A Nielsen and Molich: three to five evaluators find most problems; the CGL studies used designer panels. It is a stated limitation. Thirty-two pairs can miss small effects, which is why a non-significant result is reported as ``not detected''; the main effect, though, was large and significant.
\Q{Likert data is ordinal. Can you average it?}
\A The MOS is the ITU-T P.800 convention. The test itself uses ranks (Wilcoxon), which respects the ordinal scale; the t-test was only a sensitivity check.
\Q{Why not a separate test per format?}
\A Planned in advance: both formats tested together, each format's MOS reported without a test, so as not to multiply tests.
\Q{Isn't purposive sampling biased?}
\A The criteria ensured valid test inputs (real, unseen, text that fits). It does not claim to represent all posters; that is a limitation.
\Q{How do you know the difference came from your method?}
\A The baseline shared everything else (archetype, palette, background, markers); only the mapping and GNN differed. Because the GNN's output is registered to the anchors, the difference reflects mainly the mapping.

%=====================================================
\section{Pomelli (the panel asked for it in May)}
Pomelli is a Google Labs and Google DeepMind experiment (announced 28 October 2025). A business enters its website; Pomelli builds a ``Business DNA'' (tone, fonts, imagery, colours) and generates editable on-brand marketing images. \textbf{Answer:} ``Pomelli is industry context. It builds a brand profile from a website and generates marketing images. ThemeSync starts from the organization's own master poster, keeps each element editable, solves the layout under constraints, and keeps the designer's edits.''

%=====================================================

\needspace{14\baselineskip}
\section{Numbers to know}
\begin{tabularx}{\linewidth}{@{}Xr@{}}
\toprule
Detector (YOLOv8s, 200 images, 160/40): mAP50 / mean IoU / recall & 0.728 / 0.644 / 75\% \\
CGL-only detector on real posters (detections at 0.25) & 0 \\
GNN position error, 2,900 held-out layouts & 0.050 $\to$ 0.036 \\
GNN unmapped placement vs.\ interpolation (anchors exact) & 0.075 vs.\ 0.081 \\
Structure error, mapping + GNN vs.\ baseline: event ID / certificate & 0.045 vs.\ 0.117 / 0.056 vs.\ 0.089 \\
F1, all 64 rated layouts & 1.0 \\
Solver median / max; full edit median; full edits within 1\,s & 0.5\,ms / 2.1\,ms; 91\,ms; 100\% \\
Re-solving edits within 0.1\,s; opening a layout & 47\%; 4.6\,s \\
Stress limits: ID header / certificate signatures & 48 chars per line / overlap from 6 \\
Fitness, baseline vs.\ mapping + GNN (lower is better) & 0.0855 vs.\ 0.0879 (p = 0.030) \\
Weights $\alpha, \beta, \gamma, \delta$ & 50, 2, 1, 1.5 \\
Phase 4: experts / posters / layouts / ratings & 4 / 16 / 64 / 768 \\
MOS overall, baseline vs.\ mapping + GNN & 3.73 vs.\ 3.33 \\
Constraint and Whitespace Harmony & 3.77 vs.\ 2.92 \\
Wilcoxon overall: $p_{Holm}$ / $r$; pairs B$>$M, M$>$B, equal & $<$0.001 / $-0.95$; 25, 2, 5 \\
Kendall's $W$ overall & 0.630 (moderate) \\
Largest gap, mapping + GNN vs.\ baseline (post hoc) & 0.072 vs.\ 0.020 \\
\bottomrule
\end{tabularx}

%=====================================================

\needspace{10\baselineskip}
\section{Never say}
\begin{itemize}
\item ``Every edit takes under 0.1 s.'' (Only the solver does.)
\item ``The detector was trained on CGL.'' (COCO-pretrained; CGL trains the GNN.)
\item ``SIFT detects the logos.'' (YOLO does.)
\item ``The GNN made the layouts better.'' (It reduced structure error; the experts preferred the baseline.)
\item ``We already fixed the spacing.'' (Built after the evaluation, never rated.)
\item ``The weights were tuned to the experts' scores.'' (Calibrated on the formula's own terms, before any rating was analyzed.)
\end{itemize}

\end{document}
